% !PRINTED 84
\weberpage{84}
%<<CONT

ist, je nachdem $\lambda$ von $\mu$ verschieden ist oder nicht,
\[
 A x_\mu=0,
\]
und da nach unserer Voraussetzung $A$ von Null verschieden ist,
\[
 x_\mu=0.
\]

Wir heben den am meisten angewendeten besonderen Fall $m=n$ hervor und geben dem Satze für diesen Fall den folgenden Ausdruck:

\gesperrt{II. Wenn ein System von $n$ linearen homogenen Gleichungen mit $n$ Unbekannten eine von Null verschiedene Determinante hat, so haben sämmtliche Unbekannte den Werth Null, oder:}

\gesperrt{Wenn ein System von $n$ linearen Gleichungen mit ebenso vielen homogen vorkommenden Unbekannten eine Lösung hat, bei der nicht alle Unbekannten verschwinden, so verschwindet die Determinante des Systems.}

Unter der Determinante eines Systems von $n$ linearen homogenen Gleichungen mit $n$ Unbekannten ist hier die Determinante aus den $n^2$ Coëfficienten dieser Gleichungen verstanden.

Wir betrachten ferner den Fall, dass $\nu$ kleiner als $n$ ist. Da $m$ gleich oder grösser als $\nu$ sein muss, so wählen wir die $\nu$ ersten Gleichungen des Systems (1), und schreiben sie so:
\begin{equation}
\begin{aligned}
 a_1^{(1)}x_1+a_2^{(1)}x_2+\cdots+a_\nu^{(1)}x_\nu&=-a_{\nu+1}^{(1)}x_{\nu+1}-\cdots-a_n^{(1)}x_n,\\
 a_1^{(2)}x_1+a_2^{(2)}x_2+\cdots+a_\nu^{(2)}x_\nu&=-a_{\nu+1}^{(2)}x_{\nu+1}-\cdots-a_n^{(2)}x_n,\\
 &\vdots\\
 a_1^{(\nu)}x_1+a_2^{(\nu)}x_2+\cdots+a_\nu^{(\nu)}x_\nu&=-a_{\nu+1}^{(\nu)}x_{\nu+1}-\cdots-a_n^{(\nu)}x_n.
\end{aligned}
\tag{6}
\end{equation}
Wir bezeichnen wieder mit $\mu$ einen der Indices $1,2,\ldots,\nu$, multipliciren die Gleichungen (6) der Reihe nach mit $A_\mu^{(1)},A_\mu^{(2)},\ldots,A_\mu^{(\nu)}$ und addiren sie. Daraus folgt, wie vorhin, mit Benutzung von \S.\,22, (2), (6):
\begin{equation}
 A x_\mu=-C_{\nu+1,\mu}x_{\nu+1}-\cdots-C_{n,\mu}x_n,
\tag{7}
\end{equation}
wenn zur Abkürzung gesetzt ist:
\begin{equation}
 C_{\lambda,\mu}=\sum_{1,\nu}^{i}a_\lambda^{(i)}A_\mu^{(i)},
 \qquad \lambda=\nu+1,\nu+2,\ldots,n.
\tag{8}
\end{equation}
Nach \S.\,22, (2) ist $C_{\lambda,\mu}$ die Determinante, die aus der durch (4) definirten Determinante dadurch hervorgeht, dass man die
%CONT>>
